% ═══════════════════════════════════════════════════════════════════════
\section{La dimension de similitude}
% ═══════════════════════════════════════════════════════════════════════

\begin{frame}{Entre les courbes et les surfaces}
  \small
  \begin{itemize}
    \item Sierpiński : aire nulle, mais des segments de longueur totale
      infinie.
    \item Von Koch : longueur infinie, et aire nulle. La courbe sur un
      segment de longueur $\ell$ tient dans un triangle isocèle d'aire
      $\tfrac{\sqrt3}{12}\ell^2$ ; appliqué aux $4^n$ segments de $K_n$ :
      \[
        4^n\cdot\tfrac{\sqrt3}{12}\big(\tfrac13\big)^{2n}
        = \tfrac{\sqrt3}{12}\big(\tfrac49\big)^n \to 0.
      \]
  \end{itemize}
  \vspace{4pt}
  Une notion de dimension va situer ces objets entre les courbes et les
  surfaces, du point de vue de la mesure.
\end{frame}

\begin{frame}{Figure 3.8 --- Compter les copies}
  \FigBox[0.36\textheight]{%
    \begin{tikzpicture}[scale=1.05]
      \draw[sectioncolor, very thick] (0,0.9) -- (3,0.9);
      \foreach \x in {0,1,2,3} \draw[sectioncolor, thick] (\x,0.78) -- (\x,1.02);
      \node[etiq, align=center] at (1.5,-0.55) {$s = 3$, $N = 3 = 3^1$};
      \begin{scope}[shift={(4,0)}]
        \fill[steelblue4!22] (0,0) rectangle (2.1,2.1);
        \draw[sectioncolor, thin, step=0.7] (0,0) grid (2.1,2.1);
        \draw[sectioncolor, thick] (0,0) rectangle (2.1,2.1);
        \node[etiq, align=center] at (1.05,-0.55) {$s = 3$, $N = 9 = 3^2$};
      \end{scope}
      \begin{scope}[shift={(7.3,0)}, x={(0.7cm,0)}, y={(0,0.7cm)}, z={(0.32cm,0.24cm)}]
        \fill[steelblue4!22] (0,0,0) -- (3,0,0) -- (3,3,0) -- (0,3,0) -- cycle;
        \fill[steelblue4!34] (3,0,0) -- (3,0,3) -- (3,3,3) -- (3,3,0) -- cycle;
        \fill[steelblue4!12] (0,3,0) -- (3,3,0) -- (3,3,3) -- (0,3,3) -- cycle;
        \foreach \i in {1,2}{
          \draw[sectioncolor, thin] (\i,0,0) -- (\i,3,0) -- (\i,3,3);
          \draw[sectioncolor, thin] (0,\i,0) -- (3,\i,0) -- (3,\i,3);
          \draw[sectioncolor, thin] (3,0,\i) -- (3,3,\i) -- (0,3,\i);}
        \draw[sectioncolor, thick, line join=round] (0,0,0) -- (3,0,0) -- (3,0,3) -- (3,3,3)
          -- (0,3,3) -- (0,3,0) -- cycle (3,0,0) -- (3,3,0) -- (0,3,0) (3,3,0) -- (3,3,3);
        \node[etiq, align=center] at (1.9,-0.8,0) {$s = 3$, $N = 27 = 3^3$};
      \end{scope}
    \end{tikzpicture}}
  \small Chaque arête divisée en $s$ parties : $N=s^d$ copies de rapport
  $r=\tfrac1s$, où $d$ est la dimension usuelle :
  \[
    N = s^d \iff N = \left(\frac1r\right)^{\!d} \iff d = \frac{\ln N}{\ln(1/r)}.
  \]
  C'est le facteur $k^2$ (aires) et $k^3$ (volumes) du chapitre 1, lu à
  l'envers.
\end{frame}

\begin{frame}{Définition 3.2 --- Dimension de similitude}
  \begin{Enonce}{Définition 3.2 --- Dimension de similitude}
    Soit $E$ un ensemble autosimilaire, réunion de $N$ copies
    $w_1(E),\dots,w_N(E)$ de même rapport $r\in\intoo{0}{1}$, qui ne se
    chevauchent pas. La \defi{dimension de similitude} de $E$ est
    \[
      \dims(E) = \frac{\ln N}{\ln(1/r)},
    \]
    c'est-à-dire l'unique réel $d$ tel que $Nr^d = 1$.
  \end{Enonce}
  \small Un même ensemble se découpe de plusieurs façons, mais la
  dimension ne change pas : $N$ copies de rapport $r$ donnent, en
  redécoupant, $N^m$ copies de rapport $r^m$, et
  $\ln N^m/\ln(1/r^m) = \ln N/\ln(1/r)$. Par exemple, pour Sierpiński,
  $\ln9/\ln4 = \ln3/\ln2$.
\end{frame}

\begin{frame}{Exemple 3.5 --- Quelques dimensions de similitude}
  \begin{center}
  \small
  \begin{tabular}{lccc}
    \toprule
    Ensemble & $N$ & $r$ & $\dims(E) = \ln N/\ln(1/r)$ \\
    \midrule
    Ensemble de Cantor & $2$ & $\tfrac13$ & $\ln2/\ln3 \approx \num{0.6309}$ \\
    Segment & $3$ & $\tfrac13$ & $1$ \\
    Courbe de von Koch & $4$ & $\tfrac13$ & $\ln4/\ln3 \approx \num{1.2619}$ \\
    Triangle de Sierpiński & $3$ & $\tfrac12$ & $\ln3/\ln2 \approx \num{1.5850}$ \\
    Tapis de Sierpiński & $8$ & $\tfrac13$ & $\ln8/\ln3 \approx \num{1.8928}$ \\
    Carré & $9$ & $\tfrac13$ & $2$ \\
    \bottomrule
  \end{tabular}
  \end{center}
  \small Cantor : entre un point et un segment. Von Koch : au-delà d'une
  courbe. Tapis : presque une surface, malgré son aire nulle.
\end{frame}

\begin{frame}{Ce que mesure la dimension (raisonnement heuristique)}
  \small Supposons une « taille $d$-dimensionnelle » $M_d$ (longueur,
  aire, volume pour $d=1,2,3$) telle que $M_d\big(w(E)\big) = r^d M_d(E)$
  pour toute similitude de rapport $r$. Si $E$ est la réunion de $N$
  copies sans chevauchement :
  \[
    M_d(E) = N\,r^d\,M_d(E).
  \]
  Si $Nr^d=1$ ($d=\dims E$), toute valeur convient ; sinon $M_d(E)$ est
  nulle ou infinie. Pour von Koch :
  \begin{itemize}
    \item $d=1$ : $4\cdot\tfrac13\neq1$ ; pas nulle (deux points à
      distance $1$), donc la longueur est infinie ;
    \item $d=2$ : $4\cdot\tfrac19\neq1$ ; pas infinie ($K$ est borné),
      donc l'aire est nulle ;
    \item seule $d=\ln4/\ln3$ peut donner une valeur finie non nulle.
  \end{itemize}
  Mesurer une surface avec une longueur donne l'infini ; l'aire d'un fil
  idéal est nulle.
\end{frame}

\begin{frame}{Remarque 3.3 --- La dimension de Hausdorff ; qu'est-ce qu'une fractale ?}
  \begin{Rem}{Remarque 3.3 --- La dimension de Hausdorff}
    \footnotesize Le raisonnement reste heuristique : on n'a pas construit
    $M_d$. Hausdorff l'a fait en 1918, pour tout sous-ensemble de $\R^n$
    et tout $d\geq0$ : c'est la dimension de Hausdorff. Pour un ensemble
    autosimilaire sans chevauchement, elle coïncide avec la dimension de
    similitude (Hutchinson, 1981).
  \end{Rem}
  \small
  \begin{itemize}
    \item Mandelbrot : une fractale est un ensemble dont la dimension de
      Hausdorff dépasse strictement la dimension topologique ($0$ pour
      des points isolés, $1$ pour une courbe, $2$ pour une surface). Il
      a lui-même reconnu que cette définition exclut des cas voulus.
    \item Pour nous, une fractale : un ensemble irrégulier à toutes
      les échelles, dont la forme se reproduit (exactement ou à peu près)
      dans ses parties, et de dimension typiquement fractionnaire.
    \item Réserve : le tétraèdre de Sierpiński est une fractale de
      dimension exactement $2$.
  \end{itemize}
\end{frame}

\begin{frame}{Pour la classe 3.3 --- Logarithmes et échelles de mesure}
  \begin{Classe}{Logarithmes et échelles de mesure}
    \small
    \begin{itemize}
      \item $3^d=4$ n'a pas de solution entière : on passe aux
        logarithmes, dans la base de son choix,
        $\log4/\log3 = \ln4/\ln3 = \log_3 4$ (changement de base,
        $5^{\text{e}}$ secondaire).
      \item La dimension est un exposant : « à quelle puissance élever
        le facteur de réduction pour obtenir le nombre de copies ? »
      \item Même idée dans les échelles logarithmiques : Richter (un
        séisme de magnitude $6$ produit des ondes dix fois plus amples
        qu'un séisme de magnitude $5$, et libère environ
        $10^{\num{1.5}}\approx32$ fois plus d'énergie), pH, décibels.
      \item Calculer la magnitude de séismes historiques, puis la
        dimension de quelques fractales : deux applications d'un même
        outil.
    \end{itemize}
  \end{Classe}
\end{frame}
